\documentclass[a4paper,11pt]{article}

% Layout
\usepackage[a4paper, top=2.2cm, bottom=2.2cm, left=2.5cm, right=2.5cm]{geometry}

% Encoding / font
\usepackage[T1]{fontenc}
\usepackage[utf8]{inputenc}
\usepackage{lmodern}
\usepackage{microtype}

% Colour
\usepackage{xcolor}
\definecolor{headblue}{HTML}{1B3A6B}
\definecolor{codebg}{HTML}{F4F4F4}
\definecolor{warnbg}{HTML}{FFF3E0}
\definecolor{infobg}{HTML}{E8F4FD}
\definecolor{successbg}{HTML}{E8F5E9}

% Coloured callout boxes
\usepackage{tcolorbox}
\newtcolorbox{warnbox}{
  colback=warnbg,
  colframe=orange!70!red,
  boxrule=1pt,
  left=8pt, right=8pt, top=5pt, bottom=5pt,
  arc=2pt,
}
\newtcolorbox{infobox}{
  colback=infobg,
  colframe=blue!55!black,
  boxrule=1pt,
  left=8pt, right=8pt, top=5pt, bottom=5pt,
  arc=2pt,
}
\newtcolorbox{fixbox}{
  colback=successbg,
  colframe=green!50!black,
  boxrule=1pt,
  left=8pt, right=8pt, top=5pt, bottom=5pt,
  arc=2pt,
}

% Code listings
\usepackage{listings}
\lstset{
  basicstyle=\ttfamily\small,
  backgroundcolor=\color{codebg},
  breaklines=true,
  breakatwhitespace=false,
  postbreak=\mbox{\textcolor{gray}{$\hookrightarrow$}\ },
  frame=single,
  framerule=0.4pt,
  rulecolor=\color{gray!50},
  xleftmargin=4pt,
  xrightmargin=4pt,
  aboveskip=4pt,
  belowskip=4pt,
  columns=flexible,
  keepspaces=true,
  showstringspaces=false,
}

% Section headings
\usepackage{titlesec}
\titleformat{\section}
  {\large\bfseries\color{headblue}}{}{0em}{}
  [\vspace{1pt}\color{headblue!60}\hrule\vspace{2pt}]
\titleformat{\subsection}
  {\normalsize\bfseries\color{headblue!80}}{}{0em}{}
\titlespacing{\section}{0pt}{14pt}{5pt}
\titlespacing{\subsection}{0pt}{9pt}{3pt}

% Paragraph spacing
\usepackage{parskip}
\setlength{\parskip}{5pt}

% Lists
\usepackage{enumitem}
\setlist[enumerate]{topsep=2pt, itemsep=2pt, parsep=0pt, leftmargin=*}
\setlist[itemize]{topsep=2pt, itemsep=2pt, parsep=0pt, leftmargin=*}

% Tables
\usepackage{tabularx}
\usepackage{booktabs}
\usepackage{array}

% Header / footer
\usepackage{fancyhdr}
\pagestyle{fancy}
\fancyhf{}
\lhead{\small\color{gray!70} DashyDashboard --- Troubleshooting Guide}
\rhead{\small\color{gray!70} Broadridge Internal}
\cfoot{\small\thepage}
\renewcommand{\headrulewidth}{0.4pt}

% Links
\usepackage[hidelinks]{hyperref}

% ================================================================
\begin{document}
% ================================================================

\thispagestyle{empty}

\begin{center}
  {\LARGE\bfseries\color{headblue} DashyDashboard}\\[5pt]
  {\large\bfseries Troubleshooting Guide}\\[5pt]
  {\small ASP.NET 7 (net7.0) \quad|\quad
         IIS \quad|\quad
         Windows Authentication \quad|\quad
         SQL Server}\\[3pt]
  {\small Target SQL Server: \texttt{clipvwbpod02} \quad|\quad
         Database: \texttt{ClientsAppAttestation}}\\[3pt]
  {\small\color{gray} June 2026}
\end{center}

\vspace{4pt}
\noindent{\color{headblue}\rule{\linewidth}{1.2pt}}
\vspace{6pt}

\begin{infobox}
\textbf{How to use this guide.}
Find the symptom that matches what you see.
Each entry describes exactly what error or behaviour you observe,
why it happens, and the exact steps to fix it.
If the fix involves editing a file, the exact path and the exact text to type are given.
\end{infobox}

\tableofcontents
\newpage

% ================================================================
\section{App Refuses to Start --- ``Connection string not configured''}
\label{sec:noconnstring}
% ================================================================

\textbf{Symptom.}
The site returns a 500 error immediately.
The IIS logs or stdout log show:

\begin{lstlisting}
InvalidOperationException: Connection string 'Default' is not configured.
Set it in appsettings.Production.json or via environment variables.
\end{lstlisting}

\textbf{Why this happens.}
The app looks for the key \texttt{ConnectionStrings:Default} and throws on startup
if it is missing or blank.
This key is only in \texttt{appsettings.Production.json},
which is only loaded when \texttt{ASPNETCORE\_ENVIRONMENT=Production}.

\textbf{Fix --- check these in order:}

\begin{enumerate}
  \item Open \texttt{C:\textbackslash inetpub\textbackslash DashyDashboard\textbackslash appsettings.Production.json}
        in Notepad.
        Confirm the file contains a \texttt{ConnectionStrings} block.
        If the file is empty, missing, or the block is absent,
        paste in the full content from the Production Deployment Guide, Section~4.
  \item Confirm \texttt{ASPNETCORE\_ENVIRONMENT} is set to \texttt{Production}.
        Open \texttt{C:\textbackslash inetpub\textbackslash DashyDashboard\textbackslash web.config} in Notepad.
        Look for:
\begin{lstlisting}[language={}]
<environmentVariable name="ASPNETCORE_ENVIRONMENT"
                     value="Production" />
\end{lstlisting}
        If it is absent, add it inside the \texttt{<environmentVariables>} block
        (see the Production Deployment Guide, Section~3 for the exact structure).
  \item If you set \texttt{ASPNETCORE\_ENVIRONMENT} in IIS Manager
        (not in \texttt{web.config}), confirm it is spelled exactly
        \texttt{ASPNETCORE\_ENVIRONMENT} (all caps, underscore between words).
        A single character difference means the app loads the wrong settings file.
  \item After any change, recycle the application pool.
        Open IIS Manager, click \textbf{Application Pools},
        right-click \textbf{DashyDashboard}, click \textbf{Recycle}.
        Then browse to the site again.
\end{enumerate}

% ================================================================
\section{App Returns 500 on Every Page --- No Useful Message}
\label{sec:generic500}
% ================================================================

\textbf{Symptom.}
The browser shows a generic error page.
The response body (if any) contains only a status code and no stack trace.

\textbf{Why this happens (intentional behaviour).}
The app never leaks stack traces in production.
The global exception handler in \texttt{Program.cs} returns a JSON problem response with
no internal details.
This is by design.
You must read the actual error from the server-side logs.

\textbf{How to read the logs:}

\begin{enumerate}
  \item Open IIS Manager.
  \item In the left panel, click on the \textbf{DashyDashboard} site.
  \item In the centre panel, double-click \textbf{Logging}.
        Note the \textbf{Directory} field value (e.g.\
        \texttt{C:\textbackslash inetpub\textbackslash logs\textbackslash LogFiles}).
  \item Open File Explorer and navigate to that directory.
        Open the most recently modified \texttt{.log} file in Notepad.
        The last few lines will show the HTTP status codes and error entries.
\end{enumerate}

\textbf{To enable stdout logging for deeper errors:}

\begin{enumerate}
  \item Open \texttt{C:\textbackslash inetpub\textbackslash DashyDashboard\textbackslash web.config} in Notepad.
  \item Change \texttt{stdoutLogEnabled="false"} to \texttt{stdoutLogEnabled="true"}.
  \item Create the logs folder if it does not exist:
        open a Command Prompt as Administrator and run:
\begin{lstlisting}
mkdir C:\inetpub\DashyDashboard\logs
\end{lstlisting}
  \item Press \textbf{Ctrl + S} to save \texttt{web.config}. Browse to the site to reproduce the error.
  \item Open \texttt{C:\textbackslash inetpub\textbackslash DashyDashboard\textbackslash logs}
        and read the \texttt{stdout\_*.log} file.
  \item \textbf{Turn stdout logging off again after you are done}
        (\texttt{stdoutLogEnabled="false"}).
        Leaving it on will grow the log files indefinitely.
\end{enumerate}

% ================================================================
\section{Browser Shows ``HTTP Error 502.5 --- ANCM Out-Of-Process Startup Failure''}
\label{sec:ancm502}
% ================================================================

\textbf{Symptom.}
The IIS default error page shows:
\emph{HTTP Error 502.5 - ANCM Out-Of-Process Startup Failure}
or \emph{HTTP Error 500.30 - ASP.NET Core app failed to start}.

\textbf{Why this happens.}
The ASP.NET Core module (ANCM) started the \texttt{DashyDashboard.Api.exe} process
but the process crashed or failed to bind a port before IIS timed out.
Common causes: the .NET 7 runtime is not installed, the connection string is missing,
or the binary is for the wrong architecture.

\textbf{Fix --- work through these in order:}

\begin{enumerate}
  \item \textbf{Confirm the .NET 7 runtime is installed.}
        Open a Command Prompt on the IIS server and run:
\begin{lstlisting}
dotnet --list-runtimes
\end{lstlisting}
        You must see a line containing \texttt{Microsoft.AspNetCore.App 7.0.}
        If you do not, install the .NET 7 Hosting Bundle
        (Production Deployment Guide, Section~2)
        and then run \texttt{iisreset}.
  \item \textbf{Check the connection string.}
        A missing connection string causes an immediate crash.
        See Section~\ref{sec:noconnstring}.
  \item \textbf{Enable stdout logging and read the actual error message.}
        See Section~\ref{sec:generic500} for how to enable it.
        The stdout log will show the exact exception that killed the process.
  \item \textbf{Check the Windows Event Log.}
        Press \textbf{Win + R}, type \texttt{eventvwr}, press \textbf{Enter}.
        Expand \textbf{Windows Logs} $\rightarrow$ \textbf{Application}.
        Look for Error entries with source \textbf{IIS AspNetCore Module}
        or \textbf{DashyDashboard}.
        The event message usually contains the exact startup error.
\end{enumerate}

% ================================================================
\section{App Starts but Returns 401 Unauthorized on Every Request}
\label{sec:401}
% ================================================================

\textbf{Symptom.}
Every page returns a \texttt{401 Unauthorized} response.
The browser may show a login popup that keeps rejecting credentials,
or the app immediately shows an error page.

\textbf{Why this happens.}
Windows Authentication is not correctly configured on the IIS site.
The most common causes:
\begin{itemize}
  \item Anonymous Authentication is still enabled (it should be disabled).
  \item Windows Authentication is disabled (it should be enabled).
  \item The Negotiate provider is missing from the Windows Authentication providers list.
  \item The Windows Authentication feature was never installed on Windows Server.
\end{itemize}

\textbf{Fix:}

\begin{enumerate}
  \item Open IIS Manager.
        In the left panel, click on the \textbf{DashyDashboard} site.
        In the centre panel, double-click \textbf{Authentication}.
  \item Confirm the state of each entry:

\medskip
\begin{tabular}{@{}ll@{}}
\toprule
Authentication type & Required state \\
\midrule
Anonymous Authentication & \textbf{Disabled} \\
Windows Authentication   & \textbf{Enabled} \\
\bottomrule
\end{tabular}

\medskip
  \item If Anonymous is enabled, click it and click \textbf{Disable} in the Actions panel.
        If Windows Authentication is disabled, click it and click \textbf{Enable}.
  \item With Windows Authentication selected, click \textbf{Providers\ldots}
        Confirm \textbf{Negotiate} is in the list and is above \textbf{NTLM}.
        If Negotiate is missing, click \textbf{Add}, select \texttt{Negotiate},
        click \textbf{OK}, then click \textbf{Move Up} to put it above NTLM.
        Click \textbf{OK}.
  \item If Windows Authentication is greyed out and cannot be enabled, the Windows feature
        is not installed.
        Go to Server Manager $\rightarrow$ Add Roles and Features $\rightarrow$
        Web Server (IIS) $\rightarrow$ Web Server $\rightarrow$ Security
        and tick \textbf{Windows Authentication}. Install it.
        Then return to IIS Manager and enable it.
\end{enumerate}

% ================================================================
\section{Users Can Log In but See an Empty Dashboard (No Data)}
\label{sec:nodata}
% ================================================================

\textbf{Symptom.}
Windows Authentication works and users reach the app.
But the dashboard shows no cycles, no tools, or empty lists.
There are no JavaScript errors in the browser console.

\textbf{Why this happens.}
The database tables exist but are empty (no data was seeded),
or the migrations have not been run so the tables do not exist yet.

\textbf{Fix:}

\begin{enumerate}
  \item \textbf{Confirm the database migrations have been applied.}
        From a machine with the .NET 8 SDK that can reach \texttt{clipvwbpod02}, run:
\begin{lstlisting}
cd DashyDashboard/DashyDashboard.Api
dotnet ef migrations list --connection "Server=clipvwbpod02;Database=ClientsAppAttestation;Integrated Security=True;Encrypt=True;TrustServerCertificate=False;"
\end{lstlisting}
        Each migration in the list should show \textbf{(Applied)}.
        If any show \textbf{(Pending)}, run:
\begin{lstlisting}
dotnet ef database update --connection "Server=clipvwbpod02;Database=ClientsAppAttestation;Integrated Security=True;Encrypt=True;TrustServerCertificate=False;"
\end{lstlisting}
  \item \textbf{Confirm the app pool identity can read from the database.}
        Open SQL Server Management Studio on a machine that can reach \texttt{clipvwbpod02}.
        Run:
\begin{lstlisting}[language=SQL]
USE ClientsAppAttestation;
SELECT name FROM sys.database_principals
WHERE name = 'IIS APPPOOL\DashyDashboard';
\end{lstlisting}
        If no row is returned, the login was never created.
        Follow Step~6 in the Windows Authentication Setup section of the
        Production Deployment Guide.
  \item \textbf{Check that production data exists.}
        The app ships with no built-in production data.
        An administrator must create at least one Cycle via the admin interface
        before any user will see work to do.
        Demo data (Bruce Banner etc.) must not be used in production
        --- see Section~\ref{sec:demodata} of this guide.
\end{enumerate}

% ================================================================
\section{Users Can Log In but See ``Access Denied'' or Wrong Role}
\label{sec:wrongrole}
% ================================================================

\textbf{Symptom.}
A user who should be a Manager sees the Agent view.
Or a user who should be an Admin has no admin menu.
Or a user who should be a GFH/IFH sees nothing at all.

\textbf{Why this happens.}
Manager status is derived at runtime from whether any other user has that person's
\texttt{AssociateId} as their \texttt{ManagerId} value.
GFH/IFH/Admin roles require a row in the \texttt{SuperUsers} table.

\begin{warnbox}
\textbf{Known blocker:} the \texttt{SuperUsers} table was added in a later migration.
If the database was created from an older migration, the table may not exist.
Also, the current code stores \texttt{AssociateId} as \texttt{int},
which means any user with a non-numeric ID is silently unresolvable regardless of role.
These are tracked blockers in the codebase.
\end{warnbox}

\textbf{Fix:}

\begin{enumerate}
  \item \textbf{Check that all migrations are applied.}
        Run the \texttt{dotnet ef migrations list} command from Section~\ref{sec:nodata}.
        The migration \texttt{AddSuperUsersAndLogs} must show \textbf{(Applied)}.
        If it shows \textbf{(Pending)}, run \texttt{dotnet ef database update}.
  \item \textbf{Check the \texttt{SuperUsers} table for the affected user.}
        In SQL Server Management Studio, run:
\begin{lstlisting}[language=SQL]
SELECT * FROM dbo.SuperUsers
WHERE AssociateID = 'their-associate-id';
\end{lstlisting}
        If no row exists, the user has not been granted a SuperUser role.
        Ask the administrator to insert a row.
  \item \textbf{Check Manager status.}
        To diagnose why a user is or is not seen as a manager, run:
\begin{lstlisting}[language=SQL]
SELECT AssociateID, FirstName, LastName, ManagerId
FROM dbo.Users
WHERE ManagerId = 'the-suspected-manager-associate-id';
\end{lstlisting}
        If no rows are returned, that person has no direct reports in the database
        and will not see the Manager view.
        Ensure the direct reports' \texttt{ManagerId} column is correctly populated.
\end{enumerate}

% ================================================================
\section{SQL Server Connection Fails --- ``Login failed'' or ``Cannot open database''}
\label{sec:sqlconnfail}
% ================================================================

\textbf{Symptom.}
The app starts but every API call returns 500.
The stdout log or Windows Event Log shows an error such as:

\begin{lstlisting}
Microsoft.Data.SqlClient.SqlException: Login failed for user
  'IIS APPPOOL\DashyDashboard'.
\end{lstlisting}
or:
\begin{lstlisting}
Cannot open database "ClientsAppAttestation" requested
  by the login.
\end{lstlisting}

\textbf{Why this happens.}
The IIS application pool identity has not been granted access to the
\texttt{ClientsAppAttestation} database on \texttt{clipvwbpod02}.

\textbf{Fix:}

\begin{enumerate}
  \item Ask the DBA to run this on \texttt{clipvwbpod02}:
\begin{lstlisting}[language=SQL]
USE ClientsAppAttestation;

CREATE LOGIN [IIS APPPOOL\DashyDashboard]
    FROM WINDOWS
    WITH DEFAULT_DATABASE = ClientsAppAttestation;

CREATE USER [IIS APPPOOL\DashyDashboard]
    FOR LOGIN [IIS APPPOOL\DashyDashboard];

ALTER ROLE db_datareader
    ADD MEMBER [IIS APPPOOL\DashyDashboard];

ALTER ROLE db_datawriter
    ADD MEMBER [IIS APPPOOL\DashyDashboard];
\end{lstlisting}
  \item If the SQL runs successfully but login still fails,
        confirm the IIS application pool identity is the default
        (not a custom service account).
        Open IIS Manager $\rightarrow$ Application Pools $\rightarrow$ right-click
        \textbf{DashyDashboard} $\rightarrow$ \textbf{Advanced Settings\ldots}
        Look at the \textbf{Identity} row.
        If it is not \texttt{ApplicationPoolIdentity}, the account name in the SQL
        above is wrong.
        The account to grant access to is whatever account name appears in the
        \textbf{Identity} row.
\end{enumerate}

% ================================================================
\section{SSL / TLS Error Connecting to SQL Server}
\label{sec:sslerr}
% ================================================================

\textbf{Symptom.}
The stdout log or EF migrations output shows:

\begin{lstlisting}
A connection was successfully established with the server, but then an error occurred
during the login process. (provider: SSL Provider, error: 0 - The certificate chain
was issued by an authority that is not trusted.)
\end{lstlisting}

\textbf{Why this happens.}
The connection string contains \texttt{Encrypt=True;TrustServerCertificate=False;}.
The SQL Server on \texttt{clipvwbpod02} has a certificate that is not trusted
by the machine running the app.

\textbf{Fix options (choose one):}

\begin{enumerate}
  \item \textbf{(Preferred) Ask the DBA to install a certificate on clipvwbpod02}
        that is signed by an internal CA trusted by all servers on the network.
        This is the correct fix for a production environment.
  \item \textbf{(Temporary, for migration only) Change the migration command}
        to trust the certificate just for that one command:
        Replace \texttt{TrustServerCertificate=False} with
        \texttt{TrustServerCertificate=True} in the \texttt{--connection} argument.
        Do \textbf{not} make this change in \texttt{appsettings.Production.json}.
  \item \textbf{(Not recommended) Change the app's connection string.}
        Changing \texttt{TrustServerCertificate} to \texttt{True} in
        \texttt{appsettings.Production.json} disables certificate validation for every
        database connection the app makes.
        Do this only if the DBA confirms it is acceptable for the network environment.
\end{enumerate}

% ================================================================
\section{The Publish Command Fails with an npm Error}
\label{sec:npmfail}
% ================================================================

\textbf{Symptom.}
Running
\texttt{dotnet publish ... -c Release}
fails with an error referencing \texttt{npm} or \texttt{node}:

\begin{lstlisting}
error MSB3073: The command "npm ci --no-audit --no-fund" exited with code 1.
\end{lstlisting}
or:
\begin{lstlisting}
'npm' is not recognized as an internal or external command
\end{lstlisting}

\textbf{Why this happens.}
The publish step automatically runs \texttt{npm ci} and \texttt{npm run build}
to compile the React frontend.
Node.js and npm must be installed and on the system \texttt{PATH} on the build machine.

\textbf{Fix:}

\begin{enumerate}
  \item Open a new Command Prompt and run:
\begin{lstlisting}
node --version
npm --version
\end{lstlisting}
        If either fails, Node.js is not installed or not on the \texttt{PATH}.
        Install Node.js 18 LTS (Production Deployment Guide, Section~2).
  \item If Node is installed but the command still fails,
        the problem may be a missing \texttt{package-lock.json} file.
        Navigate to the frontend folder:
\begin{lstlisting}
cd DashyDashboard/DashyDashboard.Frontend
\end{lstlisting}
        Check whether \texttt{package-lock.json} exists:
\begin{lstlisting}
dir package-lock.json
\end{lstlisting}
        If it is missing, generate it:
\begin{lstlisting}
npm install
\end{lstlisting}
        Then go back to the repo root and run the publish command again.
  \item If \texttt{npm ci} fails with a dependency error or audit error,
        try running it manually to see the full error:
\begin{lstlisting}
cd DashyDashboard/DashyDashboard.Frontend
npm ci
\end{lstlisting}
        Read the error output and resolve the specific dependency conflict.
\end{enumerate}

% ================================================================
\section{The publish Command Fails with a dotnet SDK Version Error}
\label{sec:sdkversion}
% ================================================================

\textbf{Symptom.}
Running the publish command fails with an error such as:

\begin{lstlisting}
MSB1009: Project file does not exist.
\end{lstlisting}
or:
\begin{lstlisting}
error NU1105: Unable to find project information for ...
\end{lstlisting}
or the build exits with errors about missing targets.

\textbf{Why this happens.}
The app targets \texttt{net7.0}.
The build machine needs a .NET SDK that can compile \texttt{net7.0} targets
(SDK 7.0.x minimum; SDK 8.0.x also works).

\textbf{Fix:}

\begin{enumerate}
  \item Run this to see what SDK versions are installed:
\begin{lstlisting}
dotnet --list-sdks
\end{lstlisting}
  \item If the output is empty or shows only SDK 6.x or earlier,
        install the .NET 7 SDK.
        See the Production Deployment Guide, Section~2.
  \item If you see SDK 7.x or 8.x listed, the SDK is fine.
        The error is probably something else --- enable verbose output:
\begin{lstlisting}
dotnet build -v detailed 2>&1 | findstr /i "error"
\end{lstlisting}
        Read the first error line and resolve it specifically.
\end{enumerate}

% ================================================================
\section{dotnet ef Command Not Found}
\label{sec:efnotfound}
% ================================================================

\textbf{Symptom.}
Running \texttt{dotnet ef database update} prints:

\begin{lstlisting}
Could not execute because the specified command or file was not found.
Possible reasons for this include:
  * You intended to execute a .NET program
\end{lstlisting}

\textbf{Why this happens.}
The EF Core CLI (\texttt{dotnet-ef}) is a global tool that must be installed separately.
It is not bundled with the .NET SDK.

\textbf{Fix:}

\begin{enumerate}
  \item Install the EF Core CLI tool:
\begin{lstlisting}
dotnet tool install --global dotnet-ef
\end{lstlisting}
  \item After installation, open a \textbf{new} Command Prompt window
        (the \texttt{PATH} is updated for new windows only).
  \item Run \texttt{dotnet ef --version} to confirm it is installed.
  \item Run the migration command again.
\end{enumerate}

If you see the error again even after installing, your \texttt{PATH}
does not include the dotnet tools directory.
Open a Command Prompt and run:

\begin{lstlisting}
set PATH=%PATH%;%USERPROFILE%\.dotnet\tools
\end{lstlisting}

Then try again in the same window.

% ================================================================
\section{Demo Users Appeared in the Production Database}
\label{sec:demodata}
% ================================================================

\textbf{Symptom.}
Users with names such as \textbf{Bruce Banner}, \textbf{Diana Prince},
or \textbf{Selina Kyle} appear in the live database.

\textbf{Why this happened.}
The seed code ran in production.
This happens only when \textbf{both} of these are true at the same time:
\texttt{ASPNETCORE\_ENVIRONMENT} equals \texttt{Development},
and \texttt{DeveloperMode:EnableDemoSeedData} equals \texttt{true}.
One or both of these were true when the app started.

\textbf{Fix:}

\begin{enumerate}
  \item \textbf{Immediately} open
        \texttt{C:\textbackslash inetpub\textbackslash DashyDashboard\textbackslash appsettings.Production.json}
        in Notepad.
        Confirm all three \texttt{DeveloperMode} values are \texttt{false}.
  \item Confirm \texttt{ASPNETCORE\_ENVIRONMENT} is \texttt{Production}
        in \texttt{web.config}, not \texttt{Development}.
  \item Recycle the application pool.
  \item Remove the demo rows from the database.
        Ask the DBA to run:
\begin{lstlisting}[language=SQL]
USE ClientsAppAttestation;

-- Remove all data created by the seed (check names/IDs carefully first)
DELETE FROM dbo.ToolCycleAttestation
WHERE AssociateID IN (
    SELECT AssociateID FROM dbo.Users
    WHERE FirstName IN ('Bruce','Diana','Selina','Clark','Tony',
                        'Natasha','Peter','Steve','Wanda',
                        'James','T''Challa','Carol')
);

DELETE FROM dbo.UserToolAccess
WHERE AssociateID IN (
    SELECT AssociateID FROM dbo.Users
    WHERE FirstName IN ('Bruce','Diana','Selina','Clark','Tony',
                        'Natasha','Peter','Steve','Wanda',
                        'James','T''Challa','Carol')
);

DELETE FROM dbo.Users
WHERE FirstName IN ('Bruce','Diana','Selina','Clark','Tony',
                    'Natasha','Peter','Steve','Wanda',
                    'James','T''Challa','Carol');
\end{lstlisting}
        \textbf{Review and confirm the rows before deleting.}
        Run the \texttt{SELECT} part of each statement first
        to see exactly what will be deleted.
\end{enumerate}

% ================================================================
\section{IIS Shows the Site but the React App Does Not Load}
\label{sec:reactnotload}
% ================================================================

\textbf{Symptom.}
Navigating to the site URL shows a blank white page, a directory listing,
or the raw JSON from the API instead of the React app.

\textbf{Why this happens.}
The React frontend was not compiled into the \texttt{wwwroot} folder,
or the published output was copied incorrectly and \texttt{wwwroot/index.html}
is missing.

\textbf{Fix:}

\begin{enumerate}
  \item On the IIS server, open File Explorer and navigate to the site physical path
        (e.g.\ \texttt{C:\textbackslash inetpub\textbackslash DashyDashboard}).
  \item Look for a folder named \texttt{wwwroot}.
        Inside it there must be an \texttt{index.html} file and an \texttt{assets}
        folder containing \texttt{.js} and \texttt{.css} files.
        If \texttt{wwwroot} is empty or missing, the React build output was not
        included.
  \item Go back to the build machine and re-run the publish command
        (Production Deployment Guide, Section~7, Step~4).
        Do not skip or cancel the npm build step.
        Wait for the full \texttt{Build succeeded.} message.
  \item Copy the entire \texttt{publish\textbackslash DashyDashboard.Api} folder
        to the server again, replacing all existing files.
  \item Open a browser and navigate to the site.
        Press \textbf{Ctrl + Shift + I} to open developer tools.
        Click \textbf{Network} and press \textbf{F5} to reload.
        If you see \texttt{index.html} loading with a 200 status, the frontend
        is now being served correctly.
\end{enumerate}

% ================================================================
\section{A Specific Page Returns 404 After Reload}
\label{sec:404reload}
% ================================================================

\textbf{Symptom.}
The app loads fine from the home page.
But if a user navigates to a page
(e.g.\ \texttt{/manager} or \texttt{/attestations})
and then presses \textbf{F5} (browser refresh), they get a 404 from IIS.

\textbf{Why this happens.}
This is a single-page application (SPA).
The routes like \texttt{/manager} exist only inside React.
When the browser refreshes, it asks the server for \texttt{/manager}
as if it were a real file. The server finds no file at that path and returns 404.

\textbf{Fix:}
The app is configured to handle this automatically via:

\begin{lstlisting}[language={}]
app.MapFallbackToFile("index.html");
\end{lstlisting}

If you are still seeing 404 on refresh, the ANCM module is not routing requests
to the app. Possible causes:

\begin{enumerate}
  \item \textbf{The application pool is stopped.}
        Open IIS Manager $\rightarrow$ Application Pools.
        If \textbf{DashyDashboard} shows \textbf{Stopped}, right-click it
        and click \textbf{Start}.
  \item \textbf{Static files are being served by IIS directly, bypassing the app.}
        This can happen if a \texttt{web.config} \texttt{<staticContent>} rule is
        intercepting the request before it reaches the .NET handler.
        Remove any custom \texttt{<staticContent>} or \texttt{<handlers>} rules
        from \texttt{web.config} that you may have added.
  \item \textbf{Wrong hostingModel in web.config.}
        Confirm the \texttt{<aspNetCore>} element has
        \texttt{hostingModel="inprocess"} (not \texttt{outofprocess}).
\end{enumerate}

% ================================================================
\section{Rate Limit Error --- HTTP 429}
\label{sec:ratelimit}
% ================================================================

\textbf{Symptom.}
The app returns \texttt{HTTP 429 Too Many Requests}
for a specific user or IP address.

\textbf{Why this happens.}
The app enforces a fixed-window rate limit of 120 requests per minute,
partitioned by Windows username.
If a single user (or the same IP for unauthenticated requests) sends more than
120 requests within one 60-second window, further requests are rejected until
the window resets.

\textbf{What to check:}

\begin{itemize}
  \item This is expected behaviour, not a bug.
        A normal human user performing normal attestation work will not hit 120
        requests per minute.
  \item If an automated process or integration is hitting the limit,
        it must be slowed down.
  \item If a legitimate user is hitting the limit because of a frontend bug
        (for example, an infinite retry loop), open browser developer tools,
        click \textbf{Network}, and look for the request that is being
        sent repeatedly.
        Find and fix the offending frontend code.
\end{itemize}

% ================================================================
\section{Changes to appsettings.Production.json Are Not Being Picked Up}
\label{sec:settingsnotpicked}
% ================================================================

\textbf{Symptom.}
You edited \texttt{appsettings.Production.json} on the server but the app
behaviour has not changed.

\textbf{Why this happens.}
The app reads configuration files at startup.
Changing the file after the app has started has no effect until the app restarts.
Unlike \texttt{web.config} (which IIS monitors and reloads automatically),
\texttt{appsettings.Production.json} is only read when the process starts.

\textbf{Fix:}

\begin{enumerate}
  \item Confirm you saved the file.
        Open it in Notepad again to verify your changes are present.
  \item Recycle the application pool.
        Open IIS Manager $\rightarrow$ Application Pools $\rightarrow$
        right-click \textbf{DashyDashboard} $\rightarrow$ \textbf{Recycle}.
  \item Browse to the site.
        Your changes should now be in effect.
\end{enumerate}

\begin{infobox}
\textbf{Exception:} If you set configuration values via IIS environment variables
(in \texttt{web.config} or the Configuration Editor), those are \emph{also}
read at process start.
Saving a change to \texttt{web.config} causes IIS to restart the worker process
automatically, so environment variable changes in \texttt{web.config}
do not need a manual recycle.
\end{infobox}

% ================================================================
\section{Quick Diagnostic Checklist}
\label{sec:quickcheck}
% ================================================================

If you are unsure where a problem is, work through this checklist top to bottom.
Each item is a yes/no question.
The first ``No'' tells you where to look.

\bigskip
\renewcommand{\arraystretch}{1.5}
\begin{tabularx}{\linewidth}{|c|X|X|}
\hline
\textbf{\#} & \textbf{Check} & \textbf{If No, go to} \\
\hline
1 &
  Does \texttt{dotnet --list-runtimes} on the IIS server show
  \texttt{Microsoft.AspNetCore.App 7.0.x}? &
  Section~\ref{sec:ancm502} \\
\hline
2 &
  Does \texttt{C:\textbackslash inetpub\textbackslash DashyDashboard\textbackslash appsettings.Production.json}
  exist and contain \texttt{Server=clipvwbpod02}? &
  Section~\ref{sec:noconnstring} \\
\hline
3 &
  Is \texttt{ASPNETCORE\_ENVIRONMENT} set to \texttt{Production}
  in \texttt{web.config} or IIS? &
  Section~\ref{sec:noconnstring} \\
\hline
4 &
  Is Windows Authentication enabled and Anonymous Authentication disabled
  on the IIS site? &
  Section~\ref{sec:401} \\
\hline
5 &
  Has \texttt{IIS APPPOOL\textbackslash DashyDashboard} been granted
  db\_datareader and db\_datawriter on \texttt{ClientsAppAttestation}? &
  Section~\ref{sec:sqlconnfail} \\
\hline
6 &
  Do all migrations show \textbf{(Applied)} in
  \texttt{dotnet ef migrations list}? &
  Section~\ref{sec:nodata} \\
\hline
7 &
  Does \texttt{wwwroot/index.html} exist inside the site physical path? &
  Section~\ref{sec:reactnotload} \\
\hline
8 &
  Are all three \texttt{DeveloperMode} flags \texttt{false}
  in \texttt{appsettings.Production.json}? &
  Section~\ref{sec:demodata} \\
\hline
\end{tabularx}
\renewcommand{\arraystretch}{1}

\end{document}
